% Part: first-order-logic
% Chapter: axiomatic-proofs
% Section: identity

\documentclass[../../../include/open-logic-section]{subfiles}

\begin{document}

\olfileid{fol}{axd}{ide}

\olsection{\usetoken{S}{identity} ધરાવતી \usetoken{P}{derivation}}

!!{derivation}sમાં $\eq$ને સમાવિષ્ટ કરવા આપણે માત્ર નવાં
સ્વયંસિદ્ધિ-પ્રરૂપો ઉમેરીએ છીએ. !!{derivation} અને $\Proves$ની વ્યાખ્યા
એ જ રહે છે; હવે નવી સ્વયંસિદ્ધિઓને પણ માન્ય રાખીએ છીએ.

\begin{defn}[!!{identity} માટેની સ્વયંસિદ્ધિઓ]
\begin{align}
\ollabel{ax:id1} & \eq[t][t], \\
\ollabel{ax:id2} & \eq[t_1][t_2] \lif (!B(t_1) \lif
  !B(t_2)),
\end{align}
જ્યાં $t$, $t_1$, $t_2$ કોઈપણ બંધ પદો છે.
\end{defn}

\begin{prop}\ollabel{prop:sound}
  સ્વયંસિદ્ધિઓ \olref{ax:id1} અને \olref{ax:id2} પ્રામાણ્ય છે.
\end{prop}

\begin{proof}
સ્વાધ્યાય.
\end{proof}

\begin{prob}
\olref[fol][axd][ide]{prop:sound} સાબિત કરો.
\end{prob}

\begin{prop}\ollabel{prop:iden1}
 કોઈપણ બંધ પદ $t$ અને ગણ~$\Gamma$ માટે $\Gamma \Proves \eq[t][t]$.
\end{prop}

\begin{prop}\ollabel{prop:iden2}
  જો બંધ પદો $t_1$, $t_2$ માટે $\Gamma \Proves !A(t_1)$ અને
  $\Gamma \Proves \eq[t_1][t_2]$, તો $\Gamma \Proves !A(t_2)$.
\end{prop}

\sourcecorrection{OLAX-015}{મૂળની \texttt{identity.tex}, પંક્તિઓ
39--45માં બે પરિણામો ``કોઈપણ પદ'' માટે કહેવાય છે, જ્યારે તરત અગાઉની
\olref{ax:id1} અને \olref{ax:id2} સ્વયંસિદ્ધિઓ માત્ર બંધ પદો માટે
વ્યાખ્યાયિત છે. આપેલી સાબિતીઓ એ જ સ્વયંસિદ્ધિઓ વાપરે છે; તેથી બંને
પરિણામોમાં પદોને બંધ પદો તરીકે સ્પષ્ટ કર્યા છે.}

\begin{proof}
!!{formula}
\[
(\eq[t_1][t_2] \lif (!A(t_1) \lif !A(t_2)))
\]
એ~\olref{ax:id2}નો દૃષ્ટાંત છે. હવે~\MP{}ના બે ઉપયોગથી નિષ્કર્ષ મળે છે.
\end{proof}

\end{document}
